% Generated draft: requires mathematical and source review.
\documentclass[11pt]{article}
\usepackage[margin=25mm]{geometry}
\usepackage{fontspec}
\setmainfont{Latin Modern Roman}
\usepackage{amsmath,amssymb,mathtools,unicode-math}
\setmathfont{Latin Modern Math}
\usepackage{enumitem,needspace,xurl,hyperref}
\setlength{\parindent}{0pt}
\setlength{\parskip}{5pt}
\setlength{\emergencystretch}{3em}
\newcommand{\N}{\mathbb{N}}
\newcommand{\Z}{\mathbb{Z}}
\newcommand{\Q}{\mathbb{Q}}
\newcommand{\R}{\mathbb{R}}
\newcommand{\C}{\mathbb{C}}
\newcommand{\F}{\mathbb{F}}
\newcommand{\A}{\mathbb{A}}
\newcommand{\PP}{\mathbb{P}}
\newcommand{\E}{\mathbb{E}}
\newcommand{\eps}{\varepsilon}
\newcommand{\dd}{\,\mathrm{d}}
\begin{document}
\section*{UnsolvedMath ID: 700515}
\textbf{Source problem code:} catalogue entry
\subsection{Definitions and mathematical statement}
Use the conventional interpretation that a nil-ring is an associative ring \(R\), not necessarily with an identity, such that for every \(x\in R\) there exists an integer \(n\geq 1\) with \(x^n=0\). Its adjoint group \(R^\circ\) has underlying set \(R\) and operation
\[
x\circ y=x+y+xy,
\]
with identity \(0\). Following the supplied source, a \emph{Golod group} means a finitely generated, non-nilpotent subgroup \(G\leq R^\circ\) for some nil-ring \(R\). Here non-nilpotent means that the lower central series \(\gamma_1(G)=G\), \(\gamma_{j+1}(G)=[\gamma_j(G),G]\) never reaches the trivial subgroup.

Interpret finite rank as finite Pr\"ufer rank: an abelian group \(A\) has finite rank if there exists an integer \(r_A\geq 0\) such that every finitely generated subgroup of \(A\) can be generated by at most \(r_A\) elements.

Under these conventions, the question is whether there exist a nil-ring \(R\) and a finitely generated, non-nilpotent subgroup \(G\leq R^\circ\) such that
\[
\forall A\leq G\quad
\bigl(A\text{ abelian}\ \Longrightarrow\
\exists r_A\in\Z_{\geq 0}\ \forall H\leq A\,
(H\text{ finitely generated}\ \Longrightarrow\ d(H)\leq r_A)\bigr),
\]
where \(d(H)\) denotes the least number of generators of the finitely generated group \(H\). The bound may depend on \(A\); no common bound for all abelian subgroups is requested.

The supplied record does not explicitly define rank or the ring conventions, so the interpretation above requires confirmation against the original source. It also reports a negative answer for the whole adjoint group of a nil-ring, attributing this to Ya. P. Sysak. The exact hypotheses of that reported result are not supplied; this qualification is retained as a source report rather than expanded into an independently verified theorem.
\subsection{Short English statement}
Does some nil-ring have an adjoint group containing a finitely generated, non-nilpotent subgroup all of whose abelian subgroups have finite rank? Under the stated interpretation, rank means Pr\"ufer rank, with a separate finite bound allowed for each abelian subgroup. The supplied source reports a negative result for whole adjoint groups.
\subsection{Sources}
\begin{itemize}
\item Local UnsolvedMath registry, permanent ID \(700515\), exact source problem code \texttt{catalogue entry}, title \emph{Kourovka Notebook Problem 8.67}; supplied upstream revision \texttt{d371ee917ac9647fbb06cd6fc529deec6f3dda82}. The registry labels the problem open; that label is supplied metadata, not an independently verified mathematical assessment.
\item \emph{The Kourovka Notebook (Unsolved Problems in Group Theory)}, original problem 8.67, 8th issue (1982); proposed by A. I. Sozutov. The supplied record identifies the consulted edition as the 21st issue edition (2026), \texttt{arXiv:1401.0300v48}. The collection description names E. I. Khukhro and V. D. Mazurov as editors.
\item Ya. P. Sysak, \emph{Abstracts of the 17th All--Union Algebraic Conf., part 1}, Minsk, 1983, p. 180 (Russian). The supplied statement cites this reference for its reported negative answer concerning the whole adjoint group of a nil-ring. The reference and its precise hypotheses have not been independently checked.
\end{itemize}
\end{document}
